\documentclass[a4paper,addpoints]{article}
\usepackage{amsfonts,amssymb,amsmath, amsthm}
\usepackage{xcolor}
\usepackage{pgf,tikz,pgfplots}
\pgfplotsset{compat=1.15}
\usepgfplotslibrary{fillbetween}
\usepackage{mathrsfs}
\usetikzlibrary{arrows}
\usetikzlibrary{calc}
\usepackage{tabu}

\usepackage{tikz}
\usetikzlibrary {shapes , arrows , fit , calc, positioning}
\tikzset{box/.style={draw, rectangle, rounded corners, thick, node distance =7 em, text width =6 em, text centered, minimum height =2.5 em}}
\tikzset{line/.style={draw, thick, -latex'}}

\usepackage{tikz-uml}
\tikzumlset{actor below=0.4cm, fill usecase=white, fill class=white, text=black, draw=black}


\begin{document}



\begin{center}
\begin{tikzpicture}
    \begin{umlsystem}[x=4] {Sito per ricette}
        \umlusecase[x=4, y=3, width=1.8cm] {Cambiare lingua del sito}  
        \umlusecase[x=4, y=6, width=1.8cm] {Visualizzare homepage}
        \umlusecase[x=4, y=9, width=1.8cm] {Caso d'uso con due attori}
        \umlusecase[x=4, y=12, width=1.8cm] {Caso d'uso con due attori}
        \umlusecase[x=4, width=1.8cm] {Caso d'uso con due attori}
        \umlusecase[x=4, width=1.8cm] {Caso d'uso con due attori}
        \umlusecase[x=4, width=1.8cm] {Caso d'uso con due attori}
        \umlusecase[x=4, width=1.8cm] {Caso d'uso con due attori}
        \umlusecase[x=4, width=1.8cm] {Caso d'uso con due attori}
    \end{umlsystem}
        
    \umlactor{Visitatore} 
    
    \umlactor[x=12, y=-1.5] {Utente}

    \umlassoc{Visitatore}{usecase-1}
        
    \umlassoc{Visitatore}{usecase-4}
    \umlassoc{Utente}{usecase-4}
\end{tikzpicture}
\end{center}

\end{document} 
